% ------------------------------------------------------------
\begin{frame}{Indexes -- What They Are}
  An \textbf{index} is a separate data structure maintained by the DBMS
  alongside a table. It maps column values to the physical location of the
  corresponding row, enabling fast lookups without scanning every row.

  \medskip
  \begin{block}{Analogy}
    Think of the index at the back of a textbook. Instead of reading every
    page to find ``normalization'', you look up the word in the index and
    jump directly to the right page. A database index works the same way.
  \end{block}

  \begin{block}{Trade-offs}
    \begin{tabular}{@{}ll@{}}
      \toprule
      \textbf{Benefit}                 & \textbf{Cost} \\
      \midrule
      Fast \texttt{SELECT} / \texttt{WHERE} & Extra disk space \\
      Fast \texttt{JOIN} on foreign keys    & Slower \texttt{INSERT}/\texttt{UPDATE}/\texttt{DELETE} \\
      Fast \texttt{ORDER BY} / \texttt{GROUP BY} & Must be maintained by DBMS \\
      \bottomrule
    \end{tabular}
  \end{block}

  \begin{alertblock}{Primary and Unique keys are automatically indexed}
    MySQL always creates a \textbf{clustered B-tree index} on the primary
    key. \texttt{UNIQUE} constraints also create an index automatically.
  \end{alertblock}
\end{frame}

% ------------------------------------------------------------
\begin{frame}{B-Tree Index -- How It Works}
  The default index type in MySQL (and most relational DBMSs) is a
  \textbf{B-Tree (Balanced Tree)} index.

  \medskip
  \begin{block}{Structure}
    \begin{itemize}
      \item The tree is sorted by the indexed column value.
      \item The \textbf{root} splits the value range into subtrees.
      \item \textbf{Internal nodes} further narrow the range.
      \item \textbf{Leaf nodes} store the actual indexed values plus
            pointers to the table rows (row IDs / primary key values).
      \item The tree is always balanced -- every leaf is at the same depth,
            so lookup time is $O(\log n)$ regardless of which value is
            searched.
    \end{itemize}
  \end{block}

  \begin{block}{What B-Tree indexes support}
    \begin{itemize}
      \item \textbf{Equality}: \texttt{WHERE city = 'Berlin'} -- point lookup
      \item \textbf{Range}: \texttt{WHERE price BETWEEN 5 AND 20} -- range scan
      \item \textbf{Prefix}: \texttt{WHERE name LIKE 'Mül\%'} -- left-anchored
      \item \textbf{Sorting}: \texttt{ORDER BY} on indexed columns can skip filesort
    \end{itemize}
  \end{block}
\end{frame}

% ------------------------------------------------------------
\begin{frame}[fragile]{Creating and Dropping Indexes}
  \begin{block}{Syntax}
\begin{lstlisting}[style=sqlstyle]
-- Create a simple index on one column
CREATE INDEX idx_product_price ON product(price);

-- Create a composite (multi-column) index
CREATE INDEX idx_product_brand_price
    ON product(brand_id, price);

-- Create a unique index (also enforces uniqueness constraint)
CREATE UNIQUE INDEX uq_product_sku ON product(sku);

-- Show all indexes on a table
SHOW INDEX FROM product;

-- Remove an index
DROP INDEX idx_product_price ON product;
\end{lstlisting}
  \end{block}

  \begin{block}{When the planner uses an index}
    The planner uses the index when it estimates that an index scan is
    cheaper than a full table scan. For very small tables (a few hundred
    rows), a full scan is often faster -- the planner will choose
    accordingly.
  \end{block}
\end{frame}

% ------------------------------------------------------------
\begin{frame}{Composite Indexes and Column Order}
  A \textbf{composite index} covers multiple columns.
  The order of columns in the index matters enormously.

  \medskip
  \begin{block}{The leftmost prefix rule}
    A composite index on \texttt{(a, b, c)} can be used for:
    \begin{itemize}
      \item \texttt{WHERE a = ?} \quad -- uses index
      \item \texttt{WHERE a = ? AND b = ?} \quad -- uses index
      \item \texttt{WHERE a = ? AND b = ? AND c = ?} \quad -- uses index
      \item \texttt{WHERE b = ?} \quad -- \textbf{does NOT use index}
            (b is not the leftmost column)
      \item \texttt{WHERE a = ? AND c = ?} \quad -- uses index on \texttt{a}
            only, then filters \texttt{c} manually
    \end{itemize}
  \end{block}

  \begin{examples}{Design tip}
    Put the most selective (highest cardinality) column first.
    If queries filter by \texttt{brand\_id} then optionally by \texttt{price},
    the index should be \texttt{(brand\_id, price)}, not \texttt{(price, brand\_id)}.
  \end{examples}
\end{frame}

% ------------------------------------------------------------
\begin{frame}{Covering Indexes}
  A \textbf{covering index} is an index that contains \emph{all} the columns
  a query needs -- so the DBMS can answer the query entirely from the index
  without touching the table rows at all.

  \medskip
  \begin{block}{Example}
    Query: \texttt{SELECT price FROM product WHERE brand\_id = 5;}

    \medskip
    If the index is \texttt{(brand\_id, price)}, the engine finds the matching
    \texttt{brand\_id} entries in the index and reads \texttt{price} directly
    from the index leaf nodes. \textbf{No row lookup required.}
  \end{block}

  \begin{block}{When does it help most?}
    \begin{itemize}
      \item The table is very wide (many columns) but the query only needs a few
      \item The table has many rows and the query is executed frequently
      \item The \texttt{SELECT} list exactly matches the indexed columns
    \end{itemize}
    EXPLAIN output shows \texttt{Extra: Using index} when a covering index is used.
  \end{block}
\end{frame}

% ------------------------------------------------------------
\begin{frame}{When NOT to Use an Index}
  Indexes are not always the answer. Knowing when to \emph{avoid} them is
  equally important.

  \medskip
  \begin{block}{Situations where indexes hurt more than they help}
    \begin{itemize}
      \item \textbf{Very small tables}: a full scan of 500 rows is faster
            than an index lookup + row fetch.
      \item \textbf{Low-cardinality columns}: a \texttt{BOOLEAN} column or a
            \texttt{gender} column with only 2--3 distinct values. An index
            here hardly reduces the rows scanned.
      \item \textbf{Columns modified very frequently}: every
            \texttt{UPDATE}/\texttt{INSERT}/\texttt{DELETE} must update the
            index too. On write-heavy tables, too many indexes slow down
            writes significantly.
      \item \textbf{Functions on indexed columns}: \texttt{WHERE
            YEAR(created\_at) = 2024} defeats the index on
            \texttt{created\_at}. Rewrite as
            \texttt{WHERE created\_at >= '2024-01-01' AND created\_at < '2025-01-01'}.
    \end{itemize}
  \end{block}
\end{frame}
